\documentclass[10pt]{extarticle}
\usepackage[a4paper,margin=11mm]{geometry}
\usepackage[T1]{fontenc}
\usepackage{helvet}
\renewcommand{\familydefault}{\sfdefault}
\usepackage{amsmath,multicol,microtype,ragged2e}
\usepackage[most]{tcolorbox}
\linespread{1.2}
\setlength{\parindent}{0pt}\setlength{\parskip}{3pt}\setlength{\columnsep}{9pt}
\setlength{\abovedisplayskip}{4pt}\setlength{\belowdisplayskip}{4pt}
\newtcolorbox{qbox}[1]{enhanced,breakable,colback=white,colframe=black,boxrule=0.6pt,arc=0pt,left=5pt,right=5pt,top=3pt,bottom=3pt,before skip=6pt,after skip=6pt,
  title={\bfseries #1},coltitle=black,colbacktitle=black!12,toptitle=1pt,bottomtitle=1pt,before upper=\RaggedRight}
\pagestyle{empty}
\begin{document}
{\LARGE\bfseries Marxian Economics 101: the equations}\par\vspace{2pt}
Comrade Ammara. Examples, calculator and corrections: ahmadammar.github.io/marxist-economics\par\vspace{3pt}
\begin{multicols}{2}
\begin{qbox}{The letters}
$c$: constant capital, the value of materials and machinery used up. It passes into the product.\par
$v$: variable capital, the wages paid for labour-power.\par
$s$: surplus value, the new value beyond $v$.\par
Labour-power is the capacity to work. Labour is the activity.
\end{qbox}
\begin{qbox}{Value and cost}
\[w = c + v + s\]
\[\text{new value} = v + s\]
\[k = c + v\]
With $c=100$, $v=100$, $s=100$: $w=300$, new value $=200$, $k=200$.
\end{qbox}
\begin{qbox}{Rate of exploitation}
Surplus value for each unit of wages:
\[e = \frac{s}{v} = \frac{100}{100} = 100\%\]
\end{qbox}
\begin{qbox}{The working day}
\[T_n = T \times \frac{v}{v+s}\]
\[T_s = T - T_n\]
\[e = \frac{T_s}{T_n}\]
In an 8-hour day with $v=s$: 4 hours replace the wage, and 4 produce surplus value.
\end{qbox}
\begin{qbox}{Rate of profit}
Surplus value for each unit of capital advanced:
\[r = \frac{s}{c+v} = \frac{100}{200} = 50\%\]
\end{qbox}
\columnbreak
\begin{qbox}{Linking the rates}
Replace $s$ with $ev$, then divide top and bottom by $v$:
\[r = \frac{ev}{c+v} = \frac{e}{1+\mathrm{OCC}}\]
\[\mathrm{OCC} = \frac{c}{v}\]
If $e$ stays fixed and OCC rises, $r$ falls. Rising exploitation and cheaper machinery can offset it, so this is a tendency and not a forecast.
\end{qbox}
\begin{qbox}{Prices of production}
Competition tends to equalise profit rates:
\[r = \frac{\sum s_i}{\sum (c_i+v_i)}\]
\[p_i = (c_i+v_i)(1+r)\]
A: $50,100,100$. B: $150,100,100$. Then $r=50\%$, $p_A=225$ and $p_B=375$. Values are 250 and 350, so 25 is redistributed, and total price equals total value.
\end{qbox}
\begin{qbox}{Piece wage to time wage}
\[w_{\text{hour}} = \frac{\text{pay per piece}}{\text{hours per piece}}\]
PKR 160 per ball at 3 hours per ball gives about PKR 53.33 an hour.
\end{qbox}
\begin{qbox}{Reading the numbers with care}
Money examples assume money is proportional to value, as a teaching convention. Reported profit is only a stand-in for surplus value, so accounts cannot directly measure $e$.
\end{qbox}
\end{multicols}
\end{document}
